\chapter{A Concrete Naming Certificate for $\RapidCauchyNameUp$}
\label{ch:concrete-instances-rapid-cauchy-name-namecert}
\origin{human}

Following Bishop and Bridges, \emph{Constructive Analysis}, the $\RapidCauchyNameUp$ packet records the finite rate-controlled Cauchy real-name surface used by L10 Real and Completion consumers before any terminal $\RealUp$ seal is read. It sits over the finite window interface of \autoref{ch:concrete-instances-streamname-namecert}, the dyadic tolerance ledger of \autoref{ch:concrete-instances-dyadicratcore-namecert}, the regular rational readback boundary of \autoref{ch:concrete-instances-regseqrat-namecert}, the terminal real-name seal of \autoref{ch:concrete-instances-real-namecert}, and the Real-facing modulus route of \autoref{ch:concrete-instances-realcauchymodulus-namecert}. The packet supplies an explicit rapid Cauchy naming certificate without quotient streams, selected limits, countable choice, host real equality, trichotomy, ambient $\RealUp$ completeness, Lean verification, or a standard bridge.

\begin{definition}[Rapid Cauchy name carrier]
\label{def:rapid-cauchy-name-carrier}
A $\RapidCauchyNameUp$ carrier is a finite $\BHist$ packet
$$
R=(R_0,S,D,Q,E,H,C,P,N)
$$
with the following displayed rows:
\begin{enumerate}
\item $R_0$ is the dyadic rate schedule row, recording the finite precision requests and the displayed rule that a request at level $k$ is answered inside the dyadic tolerance budget assigned to $k$.
\item $S$ is the $\StreamNameUp$ finite-window row selected by the rate schedule for the visible requests.
\item $D$ is the $\DyadicRatCoreUp$ tolerance ledger, storing the dyadic radii, slack comparisons, and arithmetic rows consumed by the selected windows.
\item $Q$ is the $\RegSeqRatUp$ readback row obtained from the windows $S$ under the tolerance ledger $D$.
\item $E$ is the terminal $\RealUp$ seal row reached only after the rate schedule, finite windows, tolerance ledger, and regular rational readback have been displayed.
\item $H$ records componentwise $\hsame$ transport for the rate schedule, windows, tolerance ledger, readback, seal, replay, provenance, and local name.
\item $C$ records displayed $\Cont$ replay from each finite precision request through $R_0,S,D,Q,E$ and then through $H,P,N$.
\item $P$ is the $\Pkg$ provenance over $\RapidCauchyNameUp$, $\StreamNameUp$, $\DyadicRatCoreUp$, $\RegSeqRatUp$, $\RealCauchyModulusUp$, $\RealUp$, $\BHist$, $\hsame$, $\Cont$, and $\NameCert$.
\item $N$ is the local $\NameCert_{\RapidCauchyNameUp}$ row.
\end{enumerate}
The classifier compares accepted packets only through the displayed rows $R_0,S,D,Q,E,H,C,P,N$. It has no coordinate for quotient stream equality, selected limits, countable choice, host real equality, trichotomy, ambient $\RealUp$ completeness, Lean evidence, or a standard bridge.
\end{definition}
\leanchecked{BEDC.Derived.RapidCauchyNameUp.RapidCauchyNameTasteGate\_single\_carrier\_alignment}

\begin{theorem}[Rapid Cauchy name NameCert obligations]
\label{thm:rapid-cauchy-name-namecert-obligations}
For every accepted $\RapidCauchyNameUp$ carrier
$$
R=(R_0,S,D,Q,E,H,C,P,N),
$$
the local $\NameCert_{\RapidCauchyNameUp}$ surface consists exactly of the following finite obligation rows:
\begin{enumerate}
\item carrier exposure for the dyadic rate schedule $R_0$, finite $\StreamNameUp$ windows $S$, $\DyadicRatCoreUp$ tolerance ledger $D$, $\RegSeqRatUp$ readback $Q$, terminal $\RealUp$ seal $E$, and structural rows $H,C,P,N$;
\item rate admission, requiring every precision request to enter through $R_0$ before any stream window, tolerance comparison, readback, or seal row is consumed;
\item finite-window selection, requiring $S$ to contain only the windows selected by the displayed rate schedule;
\item dyadic accounting, requiring $D$ to store exactly the tolerance radii and slack comparisons consumed by the selected windows and readback rows;
\item readback and seal routing, requiring $Q$ and $E$ to remain downstream of the displayed $R_0,S,D$ route;
\item componentwise $\hsame$ transport, displayed $\Cont$ replay, $\Pkg$ provenance, and local naming through $H,C,P,N$;
\item non-escape for quotient stream equality, selected limits, countable choice, host real equality, trichotomy, ambient $\RealUp$ completeness, Lean verification, and standard-bridge data.
\end{enumerate}
No rapid Cauchy name obligation is stored outside these rows.
\end{theorem}
\begin{proof}
Project the carrier of \autoref{def:rapid-cauchy-name-carrier}. The analytic coordinates are exactly the dyadic rate schedule, finite windows, dyadic tolerance ledger, regular rational readback, and terminal seal.

The request order is forced by the tuple. A precision request first reads $R_0$, and the selected windows $S$ are visible only as windows controlled by that schedule.

The tolerance and readback rows are downstream. The ledger $D$ supplies the dyadic radii and slack comparisons consumed by $S$ and $Q$, and the seal $E$ is reached only after those rows have been displayed.

The structural rows $H,C,P,N$ transport, replay, source, and name the same finite route. Since none of the excluded objects appears as a carrier coordinate, the listed rows exhaust the local NameCert obligation surface.
\end{proof}

\begin{theorem}[Rapid Cauchy name dyadic rate route]
\label{thm:rapid-cauchy-name-dyadic-rate-route}
Every precision request accepted by a $\RapidCauchyNameUp$ carrier factors through the finite route
$$
R_0\longmapsto S\longmapsto D\longmapsto Q\longmapsto E\longmapsto H,C,P,N .
$$
Consequently the consumer receives a displayed dyadic rate schedule, selected finite $\StreamNameUp$ windows, a $\DyadicRatCoreUp$ tolerance ledger, a $\RegSeqRatUp$ readback row, a terminal $\RealUp$ seal, componentwise transport, continuation replay, provenance, and local naming. The route supplies no quotient stream equality, selected limit, countable-choice selector, host real equality, trichotomy row, ambient completeness theorem, Lean evidence, or standard bridge.
\end{theorem}
\begin{proof}
Begin with an accepted carrier and inspect the precision request. The only entrance for that request is the rate schedule row $R_0$, so the consumer cannot first read a detached stream window or a terminal Real seal.

Induct over the finite list of visible requests in the schedule. The empty list has no window or tolerance row to consume. In the cons step, the head request selects its finite window in $S$ and reads its dyadic radius and slack comparisons in $D$; the induction hypothesis handles the remaining requests.

The regular readback $Q$ is then obtained from the selected windows and tolerance ledger. The seal $E$ and structural rows $H,C,P,N$ close the same displayed route, and no carrier coordinate remains from which the excluded quotient, selector, choice, equality, order, completeness, verification, or bridge data could be read.
\end{proof}

\begin{theorem}[Rapid Cauchy name Real seal boundary]
\label{thm:rapid-cauchy-name-real-seal-boundary}
Let
$$
R=(R_0,S,D,Q,E,H,C,P,N)
$$
be an accepted $\RapidCauchyNameUp$ carrier. Any downstream use of the terminal $\RealUp$ seal $E$ must pass through the same dyadic rate schedule, selected finite windows, tolerance ledger, regular rational readback, and structural rows. Thus the seal cannot be used as a hidden endpoint for quotient stream equality, selected limits, countable choice, host real equality, trichotomy, ambient $\RealUp$ completeness, Lean verification, or a standard bridge.
\end{theorem}
\begin{proof}
Use the route of \autoref{thm:rapid-cauchy-name-dyadic-rate-route}. It places $E$ after $R_0,S,D,Q$, so every accepted seal read is attached to the displayed rate schedule, finite windows, tolerance accounting, and regular readback.

Now inspect the classifier surface of \autoref{def:rapid-cauchy-name-carrier}. The carrier has no second Real-seal coordinate and no coordinate for quotient equality, selected limits, countable choice, host equality, trichotomy, ambient completeness, Lean evidence, or bridge data.

Therefore a downstream read that bypasses one of $R_0,S,D,Q,H,C,P,N$ is not a read of the accepted carrier. The only available terminal route is the finite route named by $R$.
\end{proof}

\falsifiablePrediction{Every $\RapidCauchyNameUp$ precision read factors through a dyadic rate schedule, finite StreamName windows, a DyadicRatCore tolerance ledger, RegSeqRat readback, RealUp seal, componentwise hsame transport, displayed Cont replay, Pkg provenance, and local naming.}

\independenceWitness{streamname, dyadicratcore, regseqrat, realcauchymodulus, real: $\RapidCauchyNameUp$ is not bijective to any listed sibling because it jointly carries rate admission, finite-window selection, dyadic tolerance accounting, regular readback, terminal seal routing, transport, replay, provenance, and local naming.}

\closureat{RapidCauchyNameUp}{seedStr}

\begin{closurestatus}{\RapidCauchyNameUp}
  \constructivestory{A $\RapidCauchyNameUp$ packet is generated from a dyadic rate schedule, finite StreamName windows, DyadicRatCore tolerance rows, RegSeqRat readback, a terminal RealUp seal, componentwise hsame transport, displayed Cont replay, Pkg provenance, and a local NameCert row.}
  \theoryclosure{\seedClosure}
  \scopeclosed{\autoref{def:rapid-cauchy-name-carrier}, \autoref{thm:rapid-cauchy-name-namecert-obligations}, \autoref{thm:rapid-cauchy-name-dyadic-rate-route}, and \autoref{thm:rapid-cauchy-name-real-seal-boundary} seed the finite rapid Cauchy real-name route for L10 Real and Completion consumers.}
  \formalstatus{\unformalizedV}
  \bridgestatus{none}
  \notclaimed{This chapter does not close quotient stream equality, selected limits, countable choice, host real equality, trichotomy, ambient $\RealUp$ completeness, Lean verification, or bridge checking.}
  \upgradepath{Obligation closure requires checked carrier exposure, rate admission, window selection, dyadic accounting, readback and seal routing, dyadic rate routing, and Real-seal non-escape for BEDC.Derived.RapidCauchyNameUp.}
\end{closurestatus}
